% =====================================================================
% Gauge Neural Dynamics
% NeurIPS 2026 Workshop on Symmetry and Geometry in Neural Representations
%
% Every numeric value here comes from generated/numbers.tex and every table
% from generated/table_*.tex, both produced by scripts/make_tables.py from
% results/*/results.json. Nothing numeric is typed into this file.
% =====================================================================
\documentclass{article}

% Replace with the official neurips_<year>.sty when available; same layout.
\usepackage[preprint]{neurips_workshop}

\usepackage[utf8]{inputenc}
\usepackage[T1]{fontenc}
\usepackage{url}
\usepackage{booktabs}
\usepackage{amsfonts}
\usepackage{amsmath}
\usepackage{amssymb}
\usepackage{amsthm}
\usepackage{microtype}
\usepackage{graphicx}
\usepackage{xcolor}
\usepackage{placeins}
\usepackage[round,authoryear]{natbib}
\usepackage[colorlinks=true,linkcolor=black,citecolor=black,urlcolor=blue]{hyperref}

\graphicspath{{figures/}{../figures/}}

\theoremstyle{plain}
\newtheorem{proposition}{Proposition}

\newcommand{\GL}{\mathrm{GL}}
\newcommand{\gl}{\mathfrak{gl}}
\newcommand{\R}{\mathbb{R}}
\newcommand{\Mfd}{\mathcal{M}}
\newcommand{\Ctx}{\mathcal{C}}
\newcommand{\bch}{\mathrm{BCH}}

% Development helpers: keep the build informative when a result is not yet in.
% scripts/build_paper.sh refuses to report success while any of these fire.
\newcommand{\gndinput}[1]{\IfFileExists{#1.tex}{\input{#1}}{%
  \begin{center}\fbox{\textsf{\small [table pending]}}\end{center}}}
\newcommand{\gndfig}[2]{\IfFileExists{figures/#1}{\includegraphics[width=#2]{#1}}{%
  \fbox{\parbox[c][1.1in][c]{#2}{\centering\textsf{\small [figure pending]}}}}}

\IfFileExists{generated/numbers.tex}{% Auto-generated by scripts/make_tables.py -- do not edit.
% Every numeric value quoted in the paper text is defined here and
% comes directly from results/*/results.json.

\newcommand{\ablBchOneGcs}{--}
\newcommand{\ablDimTwenty}{--}
\newcommand{\ablDimTwo}{--}
\newcommand{\ablFullMps}{$0.887 \pm 0.000$}
\newcommand{\ablFullTransport}{$0.586 \pm 0.000$}
\newcommand{\ablNoGauge}{$0.012 \pm 0.000$}
\newcommand{\ablNoGaugeDelta}{-0.574}
\newcommand{\ablNoGroupDelta}{0.004}
\newcommand{\ablNoGroupGcs}{$0.879 \pm 0.000$}
\newcommand{\ablNoInvDelta}{--}
\newcommand{\ablNoTopoDelta}{0.002}
\newcommand{\ablNoTopoMps}{$0.866 \pm 0.000$}
\newcommand{\ablNoTransportDelta}{--}
\newcommand{\ablSeeds}{1}
\newcommand{\contCirc}{$0.471 \pm 0.024$}
\newcommand{\contCircAtPi}{$0.226 \pm 0.067$}
\newcommand{\contCircGreAtPi}{$1.39 \pm 0.35$}
\newcommand{\contHeld}{$0.688 \pm 0.002$}
\newcommand{\contHeldGre}{$0.247 \pm 0.035$}
\newcommand{\contIdentity}{$-0.024 \pm 0.095$}
\newcommand{\contLift}{$0.675 \pm 0.002$}
\newcommand{\contLiftAtPi}{$0.676 \pm 0.008$}
\newcommand{\contLiftGreAtPi}{$0.08 \pm 0.01$}
\newcommand{\contNearest}{$0.608 \pm 0.002$}
\newcommand{\contSeeds}{5}
\newcommand{\contTrained}{$0.688 \pm 0.002$}
\newcommand{\gridBestBaseGcs}{$0.992 \pm 0.000$}
\newcommand{\gridBestBaseGre}{$0.158 \pm 0.000$}
\newcommand{\gridBestBaseGreName}{PCA}
\newcommand{\gridCtrlGre}{--}
\newcommand{\gridCtrlTransport}{--}
\newcommand{\gridGLabelian}{$0.556 \pm 0.000$}
\newcommand{\gridGLgre}{$1.694 \pm 0.000$}
\newcommand{\gridGNDabelian}{$0.451 \pm 0.000$}
\newcommand{\gridGNDalg}{--}
\newcommand{\gridGNDbetti}{0\%}
\newcommand{\gridGNDcis}{$0.999 \pm 0.000$}
\newcommand{\gridGNDgcs}{$0.994 \pm 0.000$}
\newcommand{\gridGNDgre}{$0.447 \pm 0.000$}
\newcommand{\gridGNDtransport}{$0.941 \pm 0.000$}
\newcommand{\gridRotAbelian}{--}
\newcommand{\gridRotFlowGre}{--}
\newcommand{\gridRotGre}{--}
\newcommand{\gridSeeds}{1}
\newcommand{\hipBestBaseGre}{$0.213 \pm 0.010$}
\newcommand{\hipBestBaseGreName}{UMAP}
\newcommand{\hipBestBaseName}{UMAP}
\newcommand{\hipBestBaseTransport}{$0.637 \pm 0.016$}
\newcommand{\hipFlowGcs}{$0.852 \pm 0.006$}
\newcommand{\hipFlowGre}{$0.250 \pm 0.081$}
\newcommand{\hipFlowMorph}{$0.174 \pm 0.048$}
\newcommand{\hipFlowMps}{$0.734 \pm 0.008$}
\newcommand{\hipFlowTransport}{$0.731 \pm 0.013$}
\newcommand{\hipGNDalg}{--}
\newcommand{\hipGNDcis}{$0.995 \pm 0.001$}
\newcommand{\hipGNDclosure}{$0.76 \pm 0.03$}
\newcommand{\hipGNDgcs}{$0.909 \pm 0.003$}
\newcommand{\hipGNDgre}{$0.108 \pm 0.010$}
\newcommand{\hipGNDleak}{$0.067 \pm 0.023$}
\newcommand{\hipGNDmag}{$0.53 \pm 0.01$}
\newcommand{\hipGNDmorph}{$0.115 \pm 0.011$}
\newcommand{\hipGNDmps}{$0.918 \pm 0.017$}
\newcommand{\hipGNDtransport}{$0.698 \pm 0.006$}
\newcommand{\hipMAgcs}{$0.990 \pm 0.000$}
\newcommand{\hipMAmag}{$0.08 \pm 0.00$}
\newcommand{\hipPCAmorph}{$0.435 \pm 0.013$}
\newcommand{\hipSeeds}{5}
\newcommand{\hipUnpairedCis}{$0.951 \pm 0.006$}
\newcommand{\hipUnpairedTransport}{$0.664 \pm 0.006$}
\newcommand{\motBestBaseName}{Autoencoder}
\newcommand{\motBestBaseTraj}{$0.188 \pm 0.009$}
\newcommand{\motBestBaseTrajName}{ManifoldAlign}
\newcommand{\motBestBaseTransport}{$-0.103 \pm 0.138$}
\newcommand{\motCircuitEquiv}{$0.206 \pm 0.006$}
\newcommand{\motGNDcis}{$0.998 \pm 0.000$}
\newcommand{\motGNDfreqcv}{$0.332 \pm 0.025$}
\newcommand{\motGNDfreqdev}{0.076}
\newcommand{\motGNDgcs}{$0.971 \pm 0.002$}
\newcommand{\motGNDgre}{$2.457 \pm 0.296$}
\newcommand{\motGNDplane}{$48.0 \pm 4.1$}
\newcommand{\motGNDtraj}{$0.053 \pm 0.002$}
\newcommand{\motGNDtransport}{$0.876 \pm 0.008$}
\newcommand{\motMAfreqcv}{$0.066 \pm 0.007$}
\newcommand{\motMAtransport}{$-0.247 \pm 0.115$}
\newcommand{\motRefFreqcv}{$0.408 \pm 0.037$}
\newcommand{\motSeeds}{5}
}{}

\title{Gauge Neural Dynamics: Learning Context-Dependent\\Coordinate Systems in
Biological Neural Representations}

\author{%
  Anonymous Author(s)\\
  Affiliation\\
  \texttt{email}
}

\begin{document}
\maketitle

\begin{abstract}
Place fields rotate with a cue card. Grid phases shift coherently between
environments. Motor cortical trajectories for different reach directions look
like rotated copies of one another. In each case the geometry appears to be
the same object seen in a different frame, but fitting one manifold to pooled
activity has no way to say so and absorbs the difference into the residual. We
model the frame change directly: an
invariant latent state $z$ is read out through a tuning map $f$ shared across
contexts, after a context-dependent transformation $x_c = f(T_c(z))$, with $T_c$
the exponential of a context-dependent element of a Lie algebra learned from
data. Identity and exact invertibility follow from the construction; closure
under composition does not, so we measure it, using structure constants obtained
by projecting the learned commutators back onto the algebra. We define four
metrics for context-dependent geometry and apply them identically to our model
and to seven baselines. On simulations of hippocampal remapping, grid-cell phase
translation and a trained reaching network, the learned gauge recovers the
generating transformations and predicts activity in contexts held out of
training, which no baseline does. The learned algebra is abelian where the
biology says it should be and non-abelian where it should not. We also show that
a single-chart gauge field cannot cover a cyclic context variable, the failure
appearing where the topology predicts. Two conditions were built so that no
latent gauge transformation exists; the model fails on both, as intended.
\end{abstract}

% =====================================================================
\section{Introduction}
% =====================================================================

Rotate the polarising cue in a cylinder and the entire CA1 place-cell map
rotates with it \citep{muller1987effects,bostock1991experience}. Stretch the
enclosure and the map stretches \citep{okeefe1996geometric}. Morph between two
familiar shapes and it deforms in a graded way
\citep{wills2005attractor,leutgeb2005independent}. In entorhinal cortex, grid
modules shift and rotate coherently across environments
\citep{fyhn2007hippocampal} and rescale after the enclosure changes
\citep{barry2007experience}. In motor cortex, the rotational structure of the
population response differs systematically with reach condition
\citep{churchland2012neural,vyas2020computation}, and yet the geometry is stable
enough across days, and even across animals, to be recovered by alignment
\citep{gallego2020longterm,degenhart2020stabilization,safaie2023preserved}.

These observations share a form. The manifold does not appear to change shape;
it appears to be expressed in a different coordinate system. Standard latent
variable analysis cannot represent that. A single embedding is fitted to
activity pooled across conditions
\citep{cunningham2014dimensionality,gallego2017neural}, so a coherent rotation of
the map and an incoherent scattering of it are, to the model, the same amount of
unexplained variance.

Gauge theory is the natural language for a quantity defined only up to a local
choice of coordinates, together with a rule for translating between choices
\citep{kobayashi1963foundations,nakahara2003geometry}. We use it literally. Let
$z \in \Mfd$ be the latent computational state, let $f$ be the tuning of the
recorded population, and let each context $c$ act on the frame:
\begin{equation}
  x_c \;=\; f\!\left(T_c(z)\right), \qquad T_c : \Mfd \to \Mfd .
  \label{eq:gnd}
\end{equation}

The content of \eqref{eq:gnd} lies in what it forbids. Since $f$ is shared and
never sees the context, any difference between contexts must be a transformation
of the latent. Effects that are not of that form cannot be absorbed and appear
as error instead. Two of the conditions below are constructed so that no latent
gauge transformation exists, and we check that the model fails on them.

We contribute, first, a framework in which $T_c$ is the exponential of a
context-dependent element of a learned finite-dimensional Lie algebra
(Sec.~\ref{sec:method}), with identity and exact invertibility structural and
closure measured rather than assumed, together with a proof that a single-chart
field of this form cannot cover a context space on which the required group
elements wind (Prop.~\ref{prop:winding}). Second, three simulations built from
established phenomenology -- place-cell remapping, grid-cell translation on a
torus, and a reaching RNN trained by BPTT -- on which the learned transformations
recover the generating geometry (Sec.~\ref{sec:experiments}); one consequence,
that the eigenvalues of latent dynamics are context invariant while their
eigenplanes rotate, is testable in existing recordings without fitting our model.
Third, four metrics for representational consistency across contexts, each with
the diagnostic needed to keep it from being satisfied trivially
(Sec.~\ref{sec:metrics}).

% =====================================================================
\section{Related work}
% =====================================================================

Population activity is routinely summarised with latent variable models, from
linear factorisations \citep{cunningham2014dimensionality} to sequential
autoencoders \citep{pandarinath2018inferring} and contrastive embeddings
\citep{schneider2023cebra}. Context and session changes, when handled, are
handled post hoc: canonical correlation \citep{hotelling1936relations},
orthogonal Procrustes \citep{schonemann1966generalized}, or semi-supervised
manifold alignment \citep{ham2005semisupervised}, as used to stabilise
brain--computer interfaces \citep{degenhart2020stabilization} and to compare
across days and animals \citep{gallego2020longterm,safaie2023preserved}. All fit
one transformation per recorded condition, so they cannot ask whether the
transformations form a group and are undefined at a condition that was not
recorded. Both questions are central here.

The phenomena we simulate are coherent place-field remapping
\citep{muller1987effects,colgin2008understanding}, grid translation and rotation
\citep{hafting2005microstructure,fyhn2007hippocampal,stensola2012entorhinal} and
grid rescaling \citep{barry2007experience}; that a single grid module's manifold
is a torus was shown with persistent homology by \citet{gardner2022toroidal},
which our second experiment mirrors. Network models treat remapping as attractor
reorganisation \citep{burak2009accurate,low2021remapping} or as inference over a
structural code \citep{whittington2020tolman}, where we ask instead what
transformation relates the observed representations.
Equivariant networks assume a known group
\citep{cohen2016group,weiler2019general,bronstein2017geometric}, and gauge
equivariant networks handle the case where no global frame exists
\citep{cohen2019gauge,bronstein2021geometric}. Nearest to us is work that learns
a Lie algebra rather than assuming one
\citep{dehmamy2021automatic,finzi2021practical,connor2021variational,keller2021topographic};
the difference is in the target, since there the group is a prior to impose and
here it is a property of a biological system to be estimated and reported with
its defects. Our topology term follows \citet{moor2020topological}, and the
analysis uses persistent homology
\citep{edelsbrunner2002topological,carlsson2009topology,tralie2018ripser}, which
has an established role in neural coding
\citep{singh2008topological,curto2017what,chaudhuri2019intrinsic}.

% =====================================================================
\section{Method}
\label{sec:method}
% =====================================================================

\subsection{A learned Lie algebra of gauge transformations}

Let $z \in \Mfd \subseteq \R^{d}$ and let $c \in \Ctx$ be an observable context
variable: an environment label, a cue direction, a task condition. We learn
generators $G_1,\dots,G_K \in \gl(d)$ spanning $\mathfrak{g} =
\mathrm{span}\{G_k\}$ together with a context embedding $\theta : \Ctx \to
\R^{K}$, and set
\begin{equation}
  A(c) \;=\; \sum_{k=1}^{K}\theta_k(c)\,G_k,
  \qquad
  T_c(z) \;=\; \exp\!\big(A(c)\big)\, z .
  \label{eq:gauge}
\end{equation}

Three properties come for free. Smoothness of $c \mapsto T_c$ follows from that
of $\theta$ and $\exp$; invertibility is exact and in closed form,
$T_c^{-1} = \exp(-A(c))$, with no learned inverse; and $T_{c_{\mathrm{ref}}} = I$
is enforced by anchoring, $\theta(c) \mapsto \theta(c) -
\theta(c_{\mathrm{ref}})$. Anchoring adds no capacity, since \eqref{eq:gnd} is
invariant under $z \mapsto Sz$, $T_c \mapsto S T_c S^{-1}$, $f \mapsto f\circ
S^{-1}$, and anchoring merely fixes the residual left-multiplication freedom.
Generators are normalised to unit Frobenius norm, removing the $\theta
\leftrightarrow G$ scale redundancy so the defects below are comparable across
runs.

The fourth property, closure under composition, is the informative one. By
Baker--Campbell--Hausdorff, $\exp(A)\exp(B) = \exp(A + B + \tfrac12[A,B] +
\dots)$, so $\{T_c\}$ is closed if and only if $\mathfrak{g}$ is a Lie
subalgebra. We do not impose this. Projecting the commutators back onto
$\mathfrak{g}$ by least squares gives structure constants and a residual,
\begin{equation}
  [G_i, G_j] \;=\; \sum_{k} f^{k}_{ij}\, G_k \;+\; R_{ij},
  \qquad
  \delta_{\mathrm{closure}} \;=\;
  \big\langle \|R_{ij}\|_F / \|[G_i,G_j]\|_F \big\rangle_{i \neq j},
  \label{eq:closure}
\end{equation}
averaged over pairs whose commutator is not numerically zero. The $f^{k}_{ij}$
define a composition law on the algebra, $\theta_{a\cdot b} =
\bch(\theta_a,\theta_b)$, truncated at second or third order. Both the training
objective and the gauge consistency score are built on it, and it is exact when
the learned algebra closes.

Writing $g_{ab} = T_b T_a^{-1}$ makes \eqref{eq:gnd} look like a fibre bundle
over $\Ctx$, but the cocycle condition $g_{bc}g_{ab} = g_{ac}$ is then automatic,
the transitions being built from one globally defined family, and so tests
nothing. What is not automatic is whether the composite stays inside the family,
which is what \eqref{eq:closure} measures. Curvature does become measurable once
transitions are chained around a loop in context space, and we report the
resulting holonomy.

\subsection{An obstruction}

Continuous context variables expose a limitation that discrete labels hide.

\begin{proposition}[winding obstruction]
\label{prop:winding}
Let $\theta : \Ctx \to \R^{K}$ be continuous and $T_c = \exp(\sum_k \theta_k(c)
G_k) \in \GL(d,\R)^{+}$. For any loop $\gamma : S^1 \to \Ctx$, the induced loop
$t \mapsto T_{\gamma(t)}$ is null-homotopic in $\GL(d,\R)^{+}$. Hence if the
ground-truth family restricted to $\gamma$ represents a non-trivial class in
$\pi_1(\GL(d,\R)^{+}) \cong \pi_1(SO(d))$, which is $\mathbb{Z}$ for $d=2$ and
$\mathbb{Z}/2$ for $d \ge 3$, no choice of $K$ and no choice of generators can
represent it.
\end{proposition}

\begin{proof}
$\theta\circ\gamma$ is null-homotopic because $\R^K$ is contractible; let $H$
contract it to a constant. Then $\exp(\sum_k H_k(\cdot,s) G_k)$ contracts
$t \mapsto T_{\gamma(t)}$ within $\GL(d,\R)^{+}$, since $\exp$ is continuous and
lands there. A null-homotopic loop represents the trivial class.
\end{proof}

This is a statement about topology, not about optimisation or capacity, and the
remedy is the usual one: pass to a covering of $\Ctx$, or equivalently use more
than one chart. Section~\ref{sec:continuous} checks both directions.

\subsection{Non-linear gauges and shared latent dynamics}

Replacing matrices with vector fields extends \eqref{eq:gauge} to
diffeomorphisms. With $V_k(z) = G_k z + v_k + \eta\, m_k(z)$ for a small MLP
$m_k$, we take $T_c$ to be the time-one flow of $\sum_k \theta_k(c) V_k$,
integrated with RK4 and inverted by integrating the negated field. At $\eta=0$
this reduces to \eqref{eq:gauge} exactly. The bracket $[V_i,V_j] = DV_j V_i -
DV_i V_j$ is evaluated by automatic differentiation and projected pointwise onto
$\mathrm{span}\{V_k\}$, leaving \eqref{eq:closure} unchanged.

If the latent computation is shared, so is its flow. With canonical dynamics
$\dot z = F(z)$, the dynamics measured in context $c$ are the push-forward
$F_c(w) = DT_c(T_c^{-1}w)\, F(T_c^{-1}w)$, which for a linear gauge is
conjugation,
\begin{equation}
  F_c \;=\; M_c\, F\, M_c^{-1}, \qquad M_c = \exp(A(c)).
  \label{eq:conjugacy}
\end{equation}
Conjugation preserves eigenvalues, so \eqref{eq:conjugacy} predicts that the
rotation frequency of the linearised latent dynamics is shared across conditions
while the rotation plane turns with them. This can be checked in recorded data
without fitting our model, and we check it in Sec.~\ref{sec:motor}.

\subsection{Architecture and objective}

The encoder $E$ and decoder $D$ are shared MLPs that never receive the context
(Fig.~\ref{fig:architecture}, Appendix~\ref{app:details}): the encoder returns
$w_c = E(x_c)$, the gauge pulls it back to $z = T_c^{-1}(w_c)$, and the decoder
maps a frame-specific latent to activity. Giving the encoder the context would
let it absorb the transformation and leave the gauge unidentifiable. The
objective is
\begin{align}
  \mathcal{L} \;=\;
  &\underbrace{\big\|D(E(x_c)) - x_c\big\|^2}_{\text{reconstruction}}
  \;+\; \lambda_{\mathrm{t}}
  \underbrace{\big\|D\big(T_b T_a^{-1} E(x_a)\big) - x_b\big\|^2}_{\text{transport}}
  \;+\; \lambda_{\mathrm{i}}
  \underbrace{\mathbb{E}_s\,\mathrm{tr}\,\mathrm{Cov}_c\big[z_{s,c}\big]}_{\text{invariance}}
  \nonumber\\[-1pt]
  &+\; \lambda_{\mathrm{g}}
  \underbrace{\frac{\big\|T_a T_b z - T_{a\cdot b} z\big\|^2}
                   {\big\|T_a T_b z - z\big\|^2}}_{\text{group}}
  \;+\; \lambda_{\mathrm{c}}\, \delta_{\mathrm{closure}}
  \;+\; \lambda_{\mathrm{h}}\, \mathcal{L}_{H_0}(x,z)
  \;+\; \text{reg.}
  \label{eq:loss}
\end{align}
The transport term does the work: predict the response in context $b$ from one
recorded in context $a$, using only the shared decoder and the learned
transformation. The group term is normalised by how far the composite map moves
points, which stops $T \equiv \mathrm{id}$ from scoring perfectly, and
$\mathcal{L}_{H_0}$ is the persistence-signature loss of
\citet{moor2020topological}. The invariance term uses paired samples where
correspondence is available; a correspondence-free variant replaces it with a
multi-scale kernel MMD \citep{gretton2012kernel} between canonical latent
distributions. We report both.

\begin{figure}[t]
  \centering
  \gndfig{fig1_concept.pdf}{\textwidth}
  \caption{\textbf{Gauge neural dynamics.} (a) Fitting one manifold to pooled
  activity turns context structure into residual scatter. (b) Context space is
  the base of a bundle whose fibre is the latent manifold: one latent state,
  expressed in a context-dependent frame, with learned transitions between
  frames. (c) Composing the transformations of two contexts should land on the
  transformation of their composite, as predicted by BCH with the learned
  structure constants. The gap is the closure defect of Eq.~\eqref{eq:closure}.}
  \label{fig:concept}
\end{figure}

% =====================================================================
\section{Metrics}
\label{sec:metrics}
% =====================================================================

The metrics have to apply to methods that estimate no transformation at all,
otherwise the comparison is not a comparison. For such a baseline we fit an
affine map per context on the training split, take its real matrix logarithm, and
extract a $K$-dimensional algebra by PCA; every metric then runs through the same
code as for our model. This is generous, since it hands them the paired
correspondence our correspondence-free variant does not use. All values are on
held-out samples.

\textit{Gauge consistency} (GCS) compares the realised composition against the
BCH prediction, $\varepsilon(a,b) = \|T_aT_b z - T_{a\cdot b}z\| / \|T_aT_b z -
z\|$, and reports $\max(0, 1 - \langle\varepsilon\rangle)$. The denominator is
not decoration: without it a family collapsed onto the identity would score
perfectly, and one of our baselines does exactly that. We therefore always report
the transformation magnitude $\langle\|T_c(z)-z\|\rangle$ beside it.

\textit{Context invariance} (CIS) is an intraclass correlation on the canonical
latent, $\mathrm{between}/(\mathrm{between}+\mathrm{within})$ with
$\mathrm{within} = \mathbb{E}_s \mathrm{tr}\,\mathrm{Cov}_c[z_{s,c}]$ and
$\mathrm{between} = \mathrm{tr}\,\mathrm{Cov}_s[\bar z_s]$. As the context
component grows the raw ratio tends to $1/C$, not to zero, since averaging $C$
contexts shrinks noise in the per-sample mean by that factor; we therefore
chance-correct to $(\mathrm{CIS}-1/C)/(1-1/C)$, without which values would not be
comparable between five contexts and sixteen. A collapsed latent still scores
zero, having no between-sample variance.

\textit{Geometric recovery} (GRE) has to contend with the conjugation freedom. We
fit one ridge-regularised affine chart $\varphi$ from canonical latent to
ground-truth coordinate on the reference context, freeze it, and measure the
failure of $\varphi \circ T_c = T^{*}_c \circ \varphi$:
\begin{equation}
  \mathrm{GRE}_c \;=\; \big\|\varphi(T_c(z)) - T^{*}_c(u)\big\| \,\big/\,
                       \big\|T^{*}_c(u) - u\big\| .
  \label{eq:gre}
\end{equation}
Normalising by the size of the true transformation makes $\mathrm{GRE}=1$ the
score of assuming no transformation, so only values below $1$ mean anything.
Where the ground-truth coordinate is not affine in the latent a fitted chart is
unreliable; we say so rather than quoting the number, and fall back on a
chart-free alternative, the leave-one-context-out $R^2$ of predicting
$\mathrm{vec}\,\log T^{*}_c$ linearly from $\theta_c$.

\textit{Manifold preservation} (MPS) rests on manifolds related by a
diffeomorphism having the same homology, so agreement of persistence diagrams
across contexts is necessary for \eqref{eq:gnd}. We compute Vietoris--Rips
diagrams of the canonical latent per context and report $1/(1 + \langle
d_{B}\rangle)$, with $d_B$ the bottleneck distance to the diagram of the true
manifold. Bottleneck rather than Wasserstein, since the summed transport cost
grows with the number of bars. Betti numbers come from the largest multiplicative
gap in the barcode, reported with the achieved gap so a marginal call stays
visible.

\FloatBarrier
% =====================================================================
\section{Experiments}
\label{sec:experiments}
% =====================================================================

Each simulation produces paired multi-context recordings: the same latent states
seen through a shared bank of tuning curves in every context. Activity is
z-scored with statistics pooled over contexts, since per-context scaling would
remove part of what we are measuring. Models are fitted on 80\% of samples and
every number is on the remainder, as mean $\pm$ s.e.m.\ over five seeds.
Baselines are PCA, UMAP \citep{mcinnes2018umap}, an autoencoder, a VAE
\citep{kingma2014auto}, CCA, per-context PCA with orthogonal Procrustes, and
semi-supervised manifold alignment (Appendix~\ref{app:details}).

\subsection{Hippocampal remapping}
\label{sec:hippocampus}

A simulated rodent forages in a circular arena while 120 place cells tile it with
Gaussian fields. Five contexts share one bank of tuning curves and differ only in
the coordinate frame: the familiar environment, two rotations, an anisotropic
stretch-and-shear, and a radially graded twist. The last of these is a
diffeomorphism but not an affine map, and so lies outside the reach of the linear
gauge by construction.

Fig.~\ref{fig:hippocampus}a shows what has to be explained: the field of one cell
follows the cue. Fig.~\ref{fig:hippocampus}b shows the latent cloud with all
contexts overlaid; before the gauge is applied the contexts fan apart, afterwards
they collapse. Transport $R^2$ is
\hipGNDtransport{} for GND against \hipBestBaseTransport{} for the best baseline
(\hipBestBaseName), with CIS \hipGNDcis{} and GRE \hipGNDgre.

Manifold alignment deserves a comment: it attains GCS \hipMAgcs, higher than our
\hipGNDgcs, while moving the latent by only \hipMAmag{} against our \hipGNDmag.
Its transformations compose because they are almost trivial, which is the failure
mode the magnitude column exists to catch.

The morph context is where the linear gauge is provably mis-specified, being a
diffeomorphism but not an affine map. In practice the mis-specification is mild:
the linear gauge still recovers it to \hipGNDmorph{}
(Table~\ref{tab:hip-per-context}), which is worse than its \hipGNDgre{} average
but far from failure. The flow gauge, which extends the same construction to
diffeomorphisms and ought therefore to win here, does not. It predicts slightly
better (transport $R^2$ \hipFlowTransport{} against \hipGNDtransport) but
recovers the geometry worse, both overall (GRE \hipFlowGre{} against
\hipGNDgre; MPS \hipFlowMps{} against \hipGNDmps) and on the morph context
itself (\hipFlowMorph). We read this as the price of expressiveness. A gauge
flexible enough to absorb structure will absorb it, whereas the linear gauge is
forced to express the same structure as a clean transformation. Constraining the
gauge is not only a modelling convenience; it is part of what makes the
transformation identifiable at all.

The correspondence-free variant, which matches distributions instead of learning
which samples correspond, reaches transport $R^2$ \hipUnpairedTransport{} and CIS
\hipUnpairedCis; that matters, since paired samples are rarely available in real
recordings.

\begin{figure}[t]
  \centering
  \gndfig{fig3_hippocampus.pdf}{\textwidth}
  \caption{\textbf{Hippocampal remapping.} (a) Rate maps of one place cell across
  contexts; the field follows the cue coherently. (b) Latent cloud, all contexts
  overlaid, coloured by true heading. Aligned contexts give one smooth gradient;
  misaligned ones interleave. First panel: our observed-frame latent, before the
  gauge. Second: our canonical latent, after. Remainder: baselines after their own
  post-hoc alignment.}
  \label{fig:hippocampus}
\end{figure}

\gndinput{generated/table_hippocampus}

\subsection{Grid-cell geometry}
\label{sec:grid}

A single grid module's activity depends on position only through a pair of
phases, so its manifold is a torus \citep{gardner2022toroidal}. Embedding the
phase pair as $(\cos\theta_1,\sin\theta_1,\cos\theta_2,\sin\theta_2)/\sqrt2$
turns a phase translation into a pair of independent plane rotations, an element
of the maximal torus of $SO(4)$. Three predictions follow: the learned algebra
should be abelian, composition should be exact already at BCH order one, and the
canonical latent should have Betti numbers $(1,2,1)$.

Since a torus is not preserved by a general linear map, we restrict the gauge
group to isometries; this is a geometric requirement rather than a convenience,
and the unrestricted variant is worse (GRE \gridGLgre{} against \gridGNDgre,
Table~\ref{tab:grid}, Appendix~\ref{app:extra}). The predictions hold. The generators nearly commute, with
mean commutator norm \gridGNDabelian{} for unit-normalised generators, and GCS is
\gridGNDgcs. The canonical latent is a torus (Fig.~\ref{fig:grid}c,d) that
factorises into two circles carrying the two phases (Fig.~\ref{fig:grid}e,f), and
Betti numbers $(1,2,1)$ are recovered in \gridGNDbetti{} of runs.

Two further families probe the boundaries (Table~\ref{tab:grid-families}). Adding
the $60^\circ$ lattice automorphism, an integer matrix and hence a genuine torus
map, makes the group $T^2 \rtimes Z_6$ non-abelian, so a non-zero commutator is
now the right answer and the norm rises to \gridRotAbelian; being non-linear on
the $\R^4$ embedding, it defeats the linear gauge (GRE \gridRotGre) while the
flow gauge does better (\gridRotFlowGre). The third family is a control we expect
to fail: a non-integer grid rescaling \citep{barry2007experience} is well defined
on the universal cover but is not a map on the torus at all, so no latent gauge
transformation can express it. GRE there is \gridCtrlGre{} and transport $R^2$
falls to \gridCtrlTransport. The framework says in advance which phenomena lie
outside it, and the measurements agree.

\begin{figure}[t]
  \centering
  \gndfig{fig4_grid_cells.pdf}{\textwidth}
  \caption{\textbf{Grid-cell torus and translation group.} (a,b) Spatial rate maps
  of one grid cell in the reference and a phase-shifted context. (c) Canonical
  latent in its leading three principal components. (d) Its persistence barcode.
  (e,f) The latent factorises into two circles carrying the two phases. (g)
  Pairwise commutator norms of the learned generators; for pure translations the
  correct algebra is abelian. (h) Leave-one-context-out prediction of the true
  phase shift from the learned algebra coordinates.}
  \label{fig:grid}
\end{figure}

\subsection{Motor cortex}
\label{sec:motor}

A rate RNN is trained by BPTT to produce bell-shaped hand velocity toward each of
eight targets at two reach extents. We impose no symmetry on the recurrent
weights, so whatever equivariance the solution has is learned from the task. This
lets us measure a ceiling instead of assuming an exact ground truth: after
optimally aligning trial-averaged trajectories of equal-extent conditions, the
circuit's own equivariance residual is \motCircuitEquiv.

Contexts are the sixteen reach conditions. GND reaches transport $R^2$
\motGNDtransport{} while every baseline is negative, the best being
\motBestBaseTransport{} for \motBestBaseName. A negative value means the method
predicts another condition's population response worse than that condition's own
mean does. Trajectory alignment error is \motGNDtraj{} against \motBestBaseTraj{}
for the best baseline (\motBestBaseTrajName), and Fig.~\ref{fig:motor}a,b shows
the collapse directly. For the conjugacy prediction of \eqref{eq:conjugacy},
fitting linear dynamics separately per condition gives a frequency coefficient of
variation of \motGNDfreqcv{} with mean rotation-plane angle \motGNDplane{}
degrees. The circuit's own spread, measured from a PCA of the recorded activity,
is \motRefFreqcv, so the learned latent reproduces it to within
\motGNDfreqdev.

Reading that comparison needs the same care as GCS. A small frequency CV is not
in itself good: a latent that fails to separate conditions has no frequency
spread at all. Manifold alignment reaches \motMAfreqcv{} while its transport
$R^2$ is \motMAtransport. Matching the circuit's value, not undercutting it, is
what Eq.~\eqref{eq:conjugacy} predicts.

GRE is not a well-posed instrument here (\motGNDgre). Commanded hand velocity is
not an affine function of the learned latent in a context-independent way, so a
chart fitted on the reference context extrapolates unreliably onto transported
latents. That is a property of the measurement rather than of the model, and it
is why we introduced the chart-free alternative and lead with transport and
trajectory alignment for this experiment.

\begin{figure}[t]
  \centering
  \gndfig{fig5_motor_cortex.pdf}{\textwidth}
  \caption{\textbf{Motor-cortex reach conditions.} (a) Trial-averaged trajectories
  in the observed frame, one colour per reach direction. (b) The same after
  inverting the learned gauge. (c) The strongest baseline after its own
  alignment. (d) Trajectory alignment error, with and without an additional
  optimal rigid map. (e) Cross-context transport. (f) Test of
  Eq.~\eqref{eq:conjugacy}: spread of the top rotation frequency across
  conditions, with the value from recorded activity as reference.}
  \label{fig:motor}
\end{figure}

\subsection{A field over context space}
\label{sec:continuous}

Discrete labels can only ask whether a transformation was fitted per recorded
condition. The stronger question is whether the model learns a field over context
space: a map $c \mapsto T_c$ correct at cue values never observed. We present
thirteen cue orientations spanning a $120^\circ$ arc, with the context variable
being the cue direction $(\cos\alpha,\sin\alpha)$, and hold four out entirely.

The field interpolates (Fig.~\ref{fig:continuous}a,b; Table~\ref{tab:continuous}).
At held-out cue angles transport $R^2$ is \contHeld, essentially unchanged from
\contTrained{} at trained angles, against \contIdentity{} for leaving the latent
untransformed and \contNearest{} for re-using the nearest observed angle's
transformation. Per-condition alignment methods cannot be evaluated here at all,
having no transformation for a condition they never saw.

Closing the family into a full circle triggers Prop.~\ref{prop:winding}. With
context space $S^1$, mean transport $R^2$ over the circle is \contCirc, and the
failure is localised: at the antipode of the reference cue it falls to
\contCircAtPi{} with GRE \contCircGreAtPi, while near the reference it is
unimpaired (Fig.~\ref{fig:continuous}c,d). Giving the identical architecture the
universal cover, that is the unwrapped angle, removes the effect: \contLift{} on
average and \contLiftAtPi{} at $180^\circ$, with GRE \contLiftGreAtPi. The two
runs differ in nothing but the topology of the domain, which argues that this is
the predicted obstruction rather than an optimisation artefact.

\begin{figure}[t]
  \centering
  \gndfig{fig8_continuous_context.pdf}{\textwidth}
  \caption{\textbf{The gauge field over context space.} (a) Learned algebra
  coordinates as a function of cue angle, queried on a dense grid; circles mark
  cues held out of training. (b) Transport at trained and held-out cue angles,
  against the only two baselines definable for an unobserved condition. (c,d) The
  winding obstruction. On a circular context space the failure is localised at the
  branch point antipodal to the reference cue; on the universal cover it
  disappears.}
  \label{fig:continuous}
\end{figure}

\FloatBarrier
% =====================================================================
\section{Ablations and robustness}
% =====================================================================

Each ablation removes or replaces one ingredient and is refitted from scratch on
the place-cell simulation (Fig.~\ref{fig:ablations}, Table~\ref{tab:ablations},
Appendix~\ref{app:extra}). Removing the gauge, which is exactly the no-context
latent model, costs \ablNoGaugeDelta{} in transport $R^2$: the largest single
effect, and the clearest evidence that the gauge rather than the extra parameters
is doing the work. Removing the group-consistency terms costs \ablNoGroupDelta{}
in prediction while degrading GCS to \ablNoGroupGcs, so a component can matter
for geometry without mattering much for fit; removing the topology term costs
\ablNoTopoDelta{} in transport but lowers MPS from \ablFullMps{} to
\ablNoTopoMps. Truncating BCH at first order lowers GCS to \ablBchOneGcs, as it
must for a non-abelian family. The latent dimension sweep shows a useful
asymmetry: under-parameterising is costly (\ablDimTwo{} at $d=2$) while
over-parameterising is nearly free (\ablDimTwenty{} at $d=20$), and since the
true latent dimension of a recording is unknown, that is the asymmetry one
wants.

For robustness (Fig.~\ref{fig:robustness}, Appendix~\ref{app:extra}) we vary observation noise from
\robNoiseLo{} to \robNoiseHi{} times the activity standard deviation, substitute
Poisson spike counts, and vary population and sample size. Transport $R^2$ falls
from \robGNDlo{} to \robGNDhi{} across the noise range while the best baseline
reaches \robPCAhi; under Poisson counts it is \robGNDpoisson. With
\robNeuronsLo{} recorded cells it is \robGNDneuronsLo, rising to
\robGNDneuronsHi{} at \robNeuronsHi.

\FloatBarrier
% =====================================================================
\section{Discussion}
% =====================================================================

Across three simulations drawn from different phenomena, treating context as a
coordinate transformation of a shared latent rather than as extra latent variance
recovers the generating transformation and predicts activity in contexts the
model was never fitted on. The group structure is not decorative. The learned
algebra is abelian exactly where the biology says it should be, for grid-cell
phase translations, and non-abelian exactly where it should not be, once a
lattice automorphism is added; and truncating the composition law degrades
composition consistency as predicted.

One negative result is worth stating plainly. The flow gauge is a strict
generalisation of the linear one, and it fits slightly better, yet it recovers
the underlying geometry worse on every recovery measure we have -- including the
one context the linear gauge cannot represent exactly. Expressiveness and
identifiability pull in opposite directions here, and for the question this
framework is asking, the constrained gauge is the more useful instrument.

Two predictions are concrete enough to take to data. The first is
\eqref{eq:conjugacy}: if reach conditions are related by a latent gauge
transformation, the eigenvalues of the condition-wise linearised dynamics, and in
particular the rotation frequency, should be shared while the eigenplanes rotate.
This is measurable in existing motor cortex datasets without fitting anything of
ours. The second follows from Prop.~\ref{prop:winding}: a smoothly varying cyclic
context, such as a rotating cue, cannot be handled by any single continuous frame
assignment, so remapping across a full $360^\circ$ cue rotation should show a
discontinuity or hysteresis somewhere on the circle, where a non-cyclic
manipulation of comparable size should not.

These are simulations, built from published phenomenology and, for the motor
case, from a trained rather than hand-designed network; nothing here is a claim
about recorded data, and the next step is public grid-cell and motor cortex
datasets. Several limitations are intrinsic. Rate remapping, in which firing
rates change while fields stay put, is a change in $f$ and not in $T_c$, and is
outside \eqref{eq:gnd} by construction, as is global remapping modelled as a
permutation of cell identity, which is a gauge transformation in observation
space. Grid rescaling is not a map on the population manifold at all. The chart
dependence of GRE is a real weakness wherever the ground-truth coordinate is not
affine in the latent. And the linear gauge assumes a global transformation; local
deformations would need a position-dependent connection, of which the flow gauge
is only a first step. Whether real neural systems maintain gauge structure of
this kind is an empirical question, and we have tried to formulate it so that it
can be answered rather than assumed.

\label{endofbody}
\clearpage
\bibliographystyle{plainnat}
\bibliography{references}

\clearpage
\appendix
% =====================================================================
% Appendix
% =====================================================================
\section{Implementation details}
\label{app:details}

\subsection{Compute and reproducibility}

Everything runs on CPU. The models are small, under $10^6$ parameters, and the
gauge relies on \texttt{matrix\_exp} and \texttt{linalg.lstsq}, neither of which
is implemented for Apple's MPS backend. Python, NumPy and PyTorch are seeded per
run and deterministic kernels are requested. The five-seed suite takes a few
hours on eight cores. A single entry point reproduces every number and figure
from scratch; see the README.

\begin{figure}[t]
  \centering
  \gndfig{fig2_architecture.pdf}{\textwidth}
  \caption{\textbf{Architecture and objective.} Only the gauge is context
  dependent. Encoder and decoder are shared and context blind, so a context
  difference has nowhere else to go.}
  \label{fig:architecture}
\end{figure}

\subsection{Model and optimisation}

Encoder and decoder are 3-layer MLPs, 192 hidden units, layer normalisation,
GELU. The context encoder is a 2-layer $\tanh$ MLP with 64 hidden units and a
small-scale output head, followed by anchoring at the reference context. The
latent dimension is 6, except in the motor simulation where it is 8 to
accommodate the condition-independent component, and $K=6$ generators are used
throughout. AdamW, learning rate $2\times10^{-3}$, weight decay $10^{-5}$, batch
size 256, gradient clipping at norm 1, cosine annealing over 200--250 epochs.
The ablations use a shorter schedule of 120 epochs on 1500 samples, applied
identically to every configuration including the unablated model, since what is
being compared there is the effect of removing an ingredient rather than the
absolute level.
Loss weights are $\lambda_{\mathrm{t}}=1$, $\lambda_{\mathrm{i}}=1$,
$\lambda_{\mathrm{g}}=0.2$, $\lambda_{\mathrm{c}}=0.05$,
$\lambda_{\mathrm{h}}=0.5$. The gauge-dependent terms are ramped in linearly over
the first 15\% of training, so the encoder is not asked to satisfy a
transformation before it has a latent to transform. Since $\theta$ depends only
on the context, the matrix exponential is evaluated once per context per step
rather than once per sample; transport and group terms are then batched over
context pairs.

For the flow gauge, $V_k(z)=G_kz+v_k+\eta\,m_k(z)$ with $\eta=0.5$ and $m_k$ a
2-layer $\tanh$ MLP whose output layer is initialised to zero, so training starts
exactly at the linear gauge. Flows use six RK4 steps. The vector-field bracket
needs $K^2$ Jacobian-vector products, which dominates the cost, so structure
constants are cached and refreshed every 25 steps and enter the objective as
derived constants, in the spirit of a target network.

\subsection{Simulations}

\textit{Place cells.} 120 cells with Gaussian fields on a sunflower tiling of a
unit disc, width $0.22$ with log-normal jitter, log-normal peak rates. Positions
come from a momentum random walk with specular reflection at the wall, which
gives the correlated, non-uniform occupancy of real foraging. Contexts act on
position through $I$, $R(60^\circ)$, $R(-40^\circ)$ with a translation, an
anisotropic stretch-and-shear $\mathrm{diag}(1.25,0.80)$ with shear $0.22$, and a
radially graded twist $u \mapsto R(0.9\,(\|u\|/R)^2)\,u$, which is a
diffeomorphism of the disc but not affine.

\textit{Grid cells.} 100 cells in one module, spacing $0.42$, orientation
$0.21$\,rad. Rates follow the standard three-cosine construction through an
exponential non-linearity. Since $k_3 = k_2-k_1$ for a hexagonal lattice, the
module's activity depends on position only through the phase pair, which is what
makes the manifold a torus; we report the homology of a single module, matching
\citet{gardner2022toroidal}. Contexts act on the phase: translations by fixed,
well-spread offsets; a $60^\circ$ lattice rotation, whose matrix
$K R_{60}K^{-1}$ is exactly integer and hence a genuine torus automorphism; and a
non-integer rescaling by $1.37$ as the control.

\textit{Motor cortex.} A 128-unit continuous-time $\tanh$ rate RNN,
$\tau\dot x = -x + W\tanh(x) + Bu(t) + b + \sqrt{2\tau}\sigma\xi(t)$, with
$\tau=50$\,ms and $\mathrm{d}t=10$\,ms, trained by BPTT for 1500 steps to emit
hand velocity toward eight targets at two extents after a go cue, following a
300\,ms preparatory period. Reaches are slightly curved, with lateral velocity
equal to the derivative of the speed profile. The curvature is there for a
reason. With perfectly straight reaches the canonical reach-plane coordinate is
confined to a line, the second row of any readout from latent to reach plane is unconstrained
by the reference context, and no method can be scored on recovering a rotation of
that plane. Recorded activity is a random dense mixture of unit rates, as with an
electrode array; twelve trials per condition are simulated with independent
intrinsic noise. Splits are taken at the level of whole trials so no trajectory
straddles the boundary.

\subsection{Baselines}

Each baseline supplies an observed-frame latent per sample per context and an
estimated linear map from the canonical frame to each context's frame; the
canonical latent is $z_c = T_c^{-1}w_c$ and every metric applies unchanged. PCA,
UMAP, the autoencoder and the VAE embed pooled activity and have their
transformations fitted post hoc by least squares on the training split.
Procrustes runs PCA per context and aligns to the reference by orthogonal
Procrustes \citep{schonemann1966generalized}. CCA aligns each context to the
reference using the paired correspondence explicitly. Manifold alignment builds a
joint graph with within-context $k$-nearest-neighbour edges and across-context
correspondence edges, embeds with the graph Laplacian \citep{ham2005semisupervised},
and extends out of sample by random-feature ridge regression.

For the headline transport score, every method including ours is read out through
the same non-linear map: a random-Fourier-feature ridge regression
\citep{rahimi2007random} from observed-frame latent to activity, fitted on the
training split. This separates the contribution of the latent and the
transformation from decoder capacity. Where a method has a native decoder we
report its own numbers too, and they are similar.

\subsection{Metric details}

Persistence diagrams come from \texttt{ripser} \citep{tralie2018ripser} on
subsamples of up to 500 points, with filtration values normalised by the mean
pairwise distance so diagrams are comparable across embeddings of different
scale. The essential class is given the diameter of the cloud as its death time;
capping it at the largest finite death instead would make it indistinguishable
from the longest noise bar and destroy the gap that Betti estimation relies on.
Before comparing diagrams we drop bars shorter than 1\% of the longest and keep
at most the forty longest, which removes the near-diagonal noise band that would
otherwise dominate an optimal matching. Betti numbers come from the largest
multiplicative gap in the sorted lifetimes, required to exceed 1.5. That
threshold is the one genuinely tuned constant in the analysis: at the sample sizes
used here a disc produces spurious $H_1$ gaps of about 1.1 and a torus produces a
true two-bar gap of about 1.9, so 1.5 separates them, but the separation for
$H_1$ on a torus is the marginal case for any gap-based estimator and we report
the achieved gap alongside the estimate.

The post-hoc algebra fitted to a baseline's transformations uses the principal
real matrix logarithm. We record the discarded imaginary part, which is large
when a transformation lies outside the principal branch, for instance a rotation
by more than $\pi$; in that case the extracted algebra element is not unique. The
GRE chart is ridge regularised with a relative penalty, which matters when the
latent dimension exceeds that of the ground-truth coordinate. An unpenalised fit
places weights of order $1/\sigma$ on low-variance latent directions, so the
chart is accurate on the reference context but extrapolates wildly onto the
transported latents it is then applied to, penalising a method for having a rich
latent rather than for getting the transformation wrong.

\clearpage
\section{Additional results}
\label{app:extra}

\begin{figure}[t]
  \centering
  \gndfig{fig6_ablations.pdf}{\textwidth}
  \caption{\textbf{Ablations.} (a) Change in transport $R^2$ when one ingredient
  is removed. (b) The same ablations on two geometric metrics, which respond
  differently. (c) Latent dimension sweep.}
  \label{fig:ablations}
\end{figure}

\begin{figure}[t]
  \centering
  \gndfig{fig7_robustness.pdf}{\textwidth}
  \caption{\textbf{Robustness and scaling.} (a,b) Observation noise. (c) Number of
  recorded cells. (d) Number of behavioural samples.}
  \label{fig:robustness}
\end{figure}

\gndinput{generated/table_grid}

\gndinput{generated/table_motor}

\gndinput{generated/table_continuous}

\gndinput{generated/table_ablations}

\gndinput{generated/table_hippocampus_per_context}

\gndinput{generated/table_grid_families}

\gndinput{generated/table_noise}


\end{document}
